% Introducción. En APA 7 no lleva encabezado: el título repetido al inicio del texto hace de encabezado.

Este informe presenta el modelamiento como un problema de optimización con restricciones (COP), la implementación en MiniZinc \parencite{nethercote2007} y el análisis experimental del problema de las plantas de energía del Taller~2 de Programación por Restricciones. Se busca un modelo correcto para cualquier número de centrales, clientes y días y, además, una búsqueda eficiente. Las secciones siguen los apartados que pide el enunciado: descripción del problema, modelo, detalles de implementación, búsqueda, pruebas, análisis y conclusiones. Las restricciones del problema se etiquetan R1 a R5; las de dominancia, RD1 a RD4; las redundantes, RR1 y RR2; y las de ruptura de simetrías, RS1 y RS2. Las mismas etiquetas aparecen en los comentarios de \texttt{PlantaEnergia.mzn}, de modo que cada restricción del informe se puede ubicar en el modelo. El archivo \texttt{readme.txt} de la entrega enumera los archivos y explica cómo ejecutarlos.

Los experimentos de las tablas siguen el mismo protocolo. El solver principal es Gecode \parencite{gecode}, con un límite de 60~s por corrida, y los solvers de programación entera HiGHS \parencite{highs} y COIN-BC \parencite{cbc}, con brecha de optimalidad cero, dan el óptimo de referencia. El estado de cada corrida es \texttt{optimo} si el solver demostró que la solución es óptima, \texttt{sat} si encontró una solución sin demostrarlo, \texttt{insat} si demostró que no hay solución y \texttt{timeout} si no encontró ninguna en el límite. La métrica principal es el número de nodos, porque no depende del equipo. Esto importa porque los experimentos de cada sección se corrieron en equipos distintos: una corrida que termina explora exactamente los mismos nodos en todos ellos (ver Análisis), pero tarda distinto, y una corrida que llega al límite se detiene en un punto distinto en cada equipo, así que sus nodos y su mejor ganancia pueden variar. Por eso, los tiempos solo se comparan en detalle entre corridas del mismo equipo. Los datos de cada corrida están en los archivos CSV de la entrega.
